% Theory scaffold (Phase 2): conventions paragraph plus three subsections,
% per genre_conventions_theory.md §(d) Steps 0–e and §(f) sub-movements.
% Total target ~1100 words. (§2.3 fold: the original §2.3 "Gertsenshtein
% kernel" + §2.4 "Channel decomposition" merged into one §2.3 covering
% both, since each was below the BHL subsection weight.)
%
% Conventions paragraph declares signature, indices, units, M_P, kappa,
% and the Levi-Civita epsilon. Does NOT include the irreducible torsion
% decomposition labelling — that goes after the action in §2.1.

\section{Theoretical framework}%
\label{Theory}

% Conventions named paragraph (Step 0): signature, indices, units, M_P,
% kappa convention, Levi-Civita epsilon. The kappa note must reference
% \cref{sec:kappa-conventions} — the canonical disambiguation in
% docs/tex/gertsenshtein_formula.tex — to resolve the two-meanings-of-kappa
% ambiguity (issue #338). The label sec:kappa-conventions retains its
% prefix because it is defined in docs/tex/, not the manuscript scaffolding;
% see STYLE_GUIDE.md §(c).
\paragraph*{Conventions}
% TODO: declare metric signature $(+,-,-,-)$; spacetime indices $\mu,\nu,\rho$
% (Greek, coordinate) and frame indices $a,b,c$ (Latin, Lorentz);
% Levi-Civita symbol $\varepsilon^{0123} = +1$ (sign convention for the
% axial pseudovector ${}^{(3)}T$); natural units $\hbar = c = 1$;
% reduced Planck mass $\MPl$ via $\MPl^2 = (8\pi G)^{-1}$;
% the gravitational coupling $\kappa = \sqrt{2}/\MPl$, equivalently
% $\kappa^2 = 16\pi G$, used throughout for the linearised vertex coupling
% and the Gertsenshtein formula \cref{GertsenshteinFormula}; explicitly NOT
% the Einstein-action constant $\kappa_\text{E} = 8\pi G$ of some BHL
% papers (see \cref{sec:kappa-conventions} for the full disambiguation).

\subsection{Poincar\'{e} gauge theory}%
\label{PGTframework}
% Pedagogical sequence per genre_conventions_theory.md §(d) Steps (a)–(e).
% Target ~350 w.

% Step (a): Riemannian baseline (notation, not pedagogy).
\paragraph*{Riemannian baseline notation}
% TODO: one sentence on metric and Levi-Civita connection notation.

% Step (b): Promotion of connection to independent field. Palatini
% formulation declared explicitly.
\paragraph*{Independent connection in the Palatini formulation}
% TODO: one to two sentences declaring Palatini formalism (metric and
% connection independent), introducing the vierbein and spin connection.

% Step (c): Paired curvature and torsion field-strength definitions.
\paragraph*{Field strengths: curvature and torsion}
% TODO: paired adjacent equations for $R^\sigma{}_{\mu\nu\rho}$ and
% $T^\sigma{}_{\mu\nu}$ with identical syntactic structure.
\begin{align}
  % TODO: curvature definition R^sigma_{mu nu rho}
  \label{Curvature}
\end{align}
\begin{align}
  % TODO: torsion definition T^sigma_{mu nu}
  \label{Torsion}
\end{align}

% Step (d): General quadratic action with named couplings.
\paragraph*{General quadratic action}
% TODO: state that this is the most general parity-preserving quadratic
% PGT action; introduce $\Alp{n}$, $\Bet{n}$, $\Gam{n}$ couplings; gloss
% the coupling groups in one prose sentence.
\begin{align}
  % TODO: most general quadratic action
  \label{QuadraticAction}
\end{align}

% Step (e): Irreducible torsion decomposition introduced AFTER the action.
% This is also the cross-reference target for App E.
\paragraph*{Irreducible torsion decomposition}
% TODO: tensor part \T[1]{}, trace vector \T[2]{}, axial pseudovector \T[3]{}
% following Hehl 1976. Cite \cite{hehl1976spin}.
\begin{align}
  % TODO: T = T^(1) + T^(2) + T^(3) decomposition
  \label{IrreducibleTorsion}
\end{align}

\subsection{Linearisation around flat $+$ $B_0$ background}%
\label{Linearisation}
% Target ~350 w. Per genre_conventions_theory.md §(f) linearisation
% sub-movement.

\paragraph*{Background and perturbation ansatz}
% TODO: state the background — flat metric $\eta_{\mu\nu}$, uniform static
% magnetic field $\Bzero$ in the $\hat{z}$ direction. Write the perturbation
% ansatz: $g_{\mu\nu} = \eta_{\mu\nu} + h_{\mu\nu}$,
% $A_\mu = \bar{A}_\mu + \delta A_\mu$. Note: TIDAL uses g = eta + h
% (raw perturbation), NOT g = eta + kappa h. See \cref{sec:kappa-conventions}.
\begin{align}
  % TODO: perturbation ansatz
  \label{PerturbationAnsatz}
\end{align}

\paragraph*{Order-counting rule}
% TODO: state immediately after the ansatz: "We work to linear order in
% $h_{\mu\nu}$ and $\delta A_\mu$, and to first order in $\Bzero$."

% Kinetic-matrix sub-movement (third move per §(f)). Equation label is the
% cross-reference target for §4 and App E.
\paragraph*{Linearised kinetic matrix}
% TODO: "After expanding the Lagrangian to quadratic order, one finds…"
% Display the kinetic matrix without showing intermediate algebra.
\begin{align}
  % TODO: kinetic matrix (linearised quadratic Lagrangian)
  \label{KineticMatrix}
\end{align}

\subsection{Gertsenshtein conversion kernel and channel structure}%
\label{Gertsenshtein}
% Target ~400 w. Combined subsection covering both the kernel display
% (was §2.3) and the channel structure observation (was §2.4). The fold
% addresses the BHL-subsection-weight concern: each was ~250 w on its own
% which sits below the BHL norm; together they cover the natural single
% topic ("how the linearised system converts gravitons to photons, and
% which channels do so").

% Required EM-convention paragraph BEFORE the formula, per the recent
% style-intel refinement. Without it the §1 literature-tension framing
% hangs and the Palessandro–Rothman / Hwang–Noh disagreement cannot be
% navigated.
\paragraph*{Electromagnetic-field conventions}
% TODO: state the convention used ($E_i = F_{i0}$ vs the alternative);
% cite \cite{hwangnoh2023graviton} for the convention discussion. Make
% explicit that TIDAL adopts the convention used by Hwang–Noh and Dandoy
% rather than the Palessandro–Rothman alternative — this is what
% reconciles the §1 literature-tension framing.

\paragraph*{Conversion probability}
% TODO: display the formula; one prose sentence on $\Bzero$-scaling; cite
% Domcke 2301.02072 (modern derivation reference) and Ejlli 2004.02714
% (non-perturbative comparator).
\begin{align}
  \Pconv = \sin^2\!\left(\frac{\kappa\,\Bzero\,D}{2}\right)
  \label{GertsenshteinFormula}
\end{align}

\paragraph*{Validation pointer}
% TODO: one-clause cross-reference: "The quantitative reproduction of this
% kernel across our validation grid is reported in \cref{Validation},
% with the simulation-to-analytic comparison shown there." Per
% STYLE_GUIDE.md §(i), include the one-clause hint describing what is
% deferred.

% Channel structure: hx<->ax kinetic-coefficient identity (was §2.4).
% Reframed per Agent 3 review from "structural observation" to a stronger
% claim: this is a structural identity at linear order under the
% Hehl-Obukhov admissibility constraint (gr-qc/0001010), not a numerical
% coincidence. The agent recommended upgrading to "theorem" wording — we
% retain "structural identity" as a middle ground that can be sharpened
% during drafting once the algebraic argument is fully written out in
% App E.
\paragraph*{Channel decomposition}
% TODO: name the hx and ax channels and the trace channel; state that the
% hx and ax kinetic-matrix entries coincide at linear order under the
% admissibility constraint of gr-qc/0001010 — this is a structural
% identity, not a parameter-point coincidence. Cross-reference
% \cref{HxAxIdentity} (App E, where the equal-coefficient line is shown
% explicitly in the symbolic output).
